% from ega3-6-fr.tex
\oldpage[III]{168}
On a d'ailleurs
\[
  H_{q_i}(M_{\bullet\bullet\bullet}'^{(i)})
  =H_{q_i}(M_{\bullet\bullet\bullet}^{(i)}\otimes_A A')
  =\mathbf{Tor}_{q_i}^S
  (X^{(i)},S';\mathscr P_\bullet^{(i)},\mathscr O_{S'})
\]
en vertu de la définition des hypertor globaux, ce qui donne la suite $(f)$
cherchée. On peut d'autre part considérer
$M_{\bullet\bullet\bullet}'^{(i)}$ comme un bicomplexe dans lequel le premier
degré est la somme du premier et du second degré du tricomplexe
$M_{\bullet\bullet\bullet}^{(i)}$, le second degré étant le troisième degré
de ce tricomplexe; comme les modules formant les
$M_{\bullet\bullet\bullet}'^{(i)}$ sont des $A'$-modules projectifs, la théorie
générale de l'hyperhomologie montre que l'homologie du bicomplexe
\[
  M_{\bullet\bullet\bullet}'^{(1)}\otimes_{A'}
  M_{\bullet\bullet\bullet}'^{(2)}\otimes_{A'}\cdots\otimes_{A'}
  M_{\bullet\bullet\bullet}'^{(m)}\otimes_{A'} Q'_\bullet
\]
est canoniquement isomorphe à son hyperhomologie (0, 11.6.5); c'est donc
l'aboutissement d'une suite spectrale de terme $E_2$ égal à
\[
  E_{pq}^2
  =\bigoplus_{q_1+\cdots+q_{m+1}=q}
  \operatorname{Tor}_p^{A'}
  (H_{q_1}^{\mathrm{II}}(M_{\bullet\bullet\bullet}'^{(1)}),\ldots,
   H_{q_m}^{\mathrm{II}}(M_{\bullet\bullet\bullet}'^{(m)}),
   H_{q_{m+1}}^{\mathrm{II}}(Q'_\bullet)).
\]
Or, comme le second degré de $Q'_\bullet$ se réduit à $0$, on a
$H_n^{\mathrm{II}}(Q'_\bullet)=0$ pour $n\ne0$ et
$H_0^{\mathrm{II}}(Q'_\bullet)=Q'_\bullet$; la formule précédente s'écrit
aussi
\[
  E_{pq}^2
  =\bigoplus_{q_1+\cdots+q_m=q}
  \operatorname{Tor}_p^{A'}
  (H_{q_1}^{\mathrm{II}}(M_{\bullet\bullet\bullet}'^{(1)}),\ldots,
   H_{q_m}^{\mathrm{II}}(M_{\bullet\bullet\bullet}'^{(m)}),Q'_\bullet).
\]
En outre, on a
\[
  H_{q_i}^{\mathrm{II}}(M_{\bullet\bullet\bullet}'^{(i)})
  =H_{q_i}^{\mathrm{II}}(M_{\bullet\bullet\bullet}^{(i)}\otimes_A A')
  =\operatorname{Tor}_{q_i}^A(L_{\bullet\bullet}^{(i)},A')
\]
en vertu de (6.3.4); mais
$L_{-j,k}^{(i)}=C^j(\mathfrak U^{(i)},\mathscr P_k^{(i)})$, somme directe des
$\Gamma(V,\mathscr P_k^{(i)})$, où $V$ parcourt les intersections (affines) de
$j+1$ ensembles du recouvrement $\mathfrak U^{(i)}$; si
$V'=u_i^{-1}(V)$, $V'$ est affine dans $X'^{(i)}$ et il résulte de (6.4.1.1)
que l'on a
\[
  \Gamma\bigl(V',\mathscr{T}\!or_{q_i}^S
  (\mathscr P_k^{(i)},\mathscr O_{S'})\bigr)
  =\operatorname{Tor}_{q_i}^A(\Gamma(V,\mathscr P_k^{(i)}),A')
\]
d'où pour le bicomplexe
$H_{q_i}^{\mathrm{II}}(M_{\bullet\bullet\bullet}'^{(i)})$ l'expression
\[
  C^\bullet\bigl(\mathfrak U'^{(i)},
  \mathscr{T}\!or_{q_i}^S(\mathscr P_\bullet^{(i)},\mathscr O_{S'})\bigr),
\]
ce qui donne finalement l'expression cherchée pour le terme $E_2$ de la
suite $(f')$. Le fait que cette suite soit birégulière sous les conditions
indiquées se vérifie comme d'ordinaire, tenant compte de ce que, si
$\mathscr P_\bullet^{(i)}$ est limité inférieurement, tous les degrés de
$M_{\bullet\bullet\bullet}'^{(i)}$ sont limités inférieurement.
\end{proof}

\begin{remark}[6.9.5]
On voit comme dans (6.7.6) que le remplacement des
$\mathscr P_\bullet^{(i)}$ et de $\mathscr Q'_\bullet$ par des complexes qui
leur sont respectivement homotopes ne change pas les suites $(e)$, $(f)$ et
$(f')$ à un isomorphisme canonique près. En outre, pour la suite $(f)$, des
homomorphismes
$\mathscr P_\bullet^{(i)}\to\mathscr R_\bullet^{(i)}$,
$\mathscr Q'_\bullet\to\mathscr T'_\bullet$ de complexes qui donnent des
isomorphismes en homologie
$\mathscr H_\bullet(\mathscr P_\bullet^{(i)})\xrightarrow{\sim}
\mathscr H_\bullet(\mathscr R_\bullet^{(i)})$,
$\mathscr H_\bullet(\mathscr Q'_\bullet)\xrightarrow{\sim}
\mathscr H_\bullet(\mathscr T'_\bullet)$ fournissent un isomorphisme de suites
spectrales
\[
\begin{split}
  {}^{(f)}\mathscr E
  (f_1,\ldots,f_m;
   \mathscr P_\bullet^{(1)},\ldots,\mathscr P_\bullet^{(m)},
   \mathscr Q'_\bullet)
  \xrightarrow{\sim}\;&
  {}^{(f)}\mathscr E
  (f_1,\ldots,f_m;\\
  &\mathscr R_\bullet^{(1)},\ldots,\mathscr R_\bullet^{(m)},
   \mathscr T'_\bullet);
\end{split}
\]
la démonstration est la même que pour (6.7.6) en tenant compte du résultat
de (6.7.6) et de la régularité de la suite $(f)$.
\end{remark}

\begin{corollary}[6.9.6]
Sous les conditions de (6.9.1), supposons que :
\begin{enumerate}
  \item[$1^\circ$] Les complexes $\mathscr P_\bullet^{(i)}$ sont formés de
  Modules plats sur $S$, et les $\mathscr O_Y$-Modules
  \[
    \mathfrak{Tor}_n^S
    (f_1,\ldots,f_m;
     \mathscr P_\bullet^{(1)},\ldots,\mathscr P_\bullet^{(m)})
  \]
  sont plats sur $S$.
  \item[$2^\circ$] Les $\mathscr P_\bullet^{(i)}$ et
  $\mathscr Q'_\bullet$ sont limités inférieurement.
\end{enumerate}

